% Cập nhật ngày 15/10/2026; S15-01--S15-04 và các nguồn đã ghi trong bản trước.
\chapter{Kết luận}
\label{chuong:ketluan}

Khóa luận đã đặc tả, cài đặt và đánh giá trên mô phỏng một trần cho tổng phần dời giờ trả do điều phối của từng
hành khách. Phương pháp kết hợp giới hạn phần điều phối kiểm soát được, lời hứa có tính trước hồ sơ tắc đường
và sổ theo từng khách. Chương này tổng hợp các kết quả, những đánh đổi quan sát được và phần còn cần kiểm tra.

\section{Kết quả và đánh đổi}
\label{muc:c5-tongket}

Luật C cộng độ lớn của từng lần dời giờ trả so với mốc giữ nguyên kế hoạch đã được tính lại theo thời gian đi
hiện tại. Mốc này kế thừa công thức 7 của Hörl và Zwick~\cite{horl2022}. Dưới điều kiện giữ nguyên kế hoạch ở
mục~\ref{muc:giunguyen}, phương án giữ nguyên không tiêu thêm trần. Trên tập kiểm tra chính, không khách nào
vượt trần 30 giây ở cả hai thứ tự chọn phương án; tổng dời lớn nhất là 29,239 giây khi ưu tiên ổn định và
29,657 giây khi ưu tiên chi phí kế hoạch. Đây là giới hạn phần tính cho điều phối, không phải bảo đảm giờ đến
thực tế cách lời hứa đầu nhiều nhất 30 giây.
% EV: A01; A06; A13; A22; X3-27; J01; J03; S15-01 (5d); S15-02.

Với ưu tiên ổn định, thêm trần 30 giây làm số lượt phục vụ giảm từ 2.434 xuống 2.338 và số khách hoàn thành
có đi chung giảm từ 555 xuống 313. Đổi lại, số lượt được phục vụ và không muộn quá một phút tăng từ 2.231
lên 2.262 trên 3.672 yêu cầu. Với ưu tiên chi phí, các số phục vụ là 2.502 và 2.364, số khách hoàn thành có
đi chung là 1.027 và 432, còn số lượt không muộn quá một phút là 2.123 và 2.292. Trần vẫn đi cùng độ trễ
thấp hơn khi bỏ ưu tiên ổn định khỏi thứ tự xếp hạng, nhưng đánh đổi về phục vụ và ghép khách cũng lớn hơn.
% EV: S15-01 (mục 2, 5f); S15-02. Nhóm khách phục vụ thay đổi giữa các cấu hình.

Khi cùng giữ trần, ưu tiên chi phí phục vụ thêm 26 lượt và có thêm 119 khách hoàn thành có đi chung so với
ưu tiên ổn định. Cái giá nằm ở nhiều khách bị sửa lời hứa hơn: 106 thay vì 33; số khách bị sửa từ hai lần
là 2 thay vì 0. Số lượt đón muộn là 44 thay vì 43. Vì thế, kết quả không thay thế bộ chọn ưu tiên ổn định
bằng một bộ chọn tốt hơn ở mọi tiêu chí. Hai thứ tự cho hai cách sử dụng phần linh hoạt còn lại trong cùng
trần; lựa chọn cần dựa trên mục tiêu dịch vụ và dữ liệu về mức chấp nhận của khách.
% EV: S15-01 (5d--5f); S15-04.

Các đối chứng cho thấy lợi ích có giới hạn. Dưới ưu tiên ổn định, luật hạn chót 30 giây cũng phục vụ 2.338
lượt; khi so cùng ngưỡng muộn hơn lời hứa đầu quá 30 giây, luật C có 292 và luật hạn chót có 289 lượt muộn.
Khi có dự báo, điều phối không trần đệm thêm 30 giây vào mọi lời hứa cũng đạt 2.328 lượt không muộn quá một
phút, bằng luật C, và hơn ở hai phút: 2.366 so với 2.338. Đổi lại, phần lời hứa dài thêm trung bình là 37,2
thay vì 7,4 giây. Những kết quả này không cho phép chọn luật C chỉ từ một số đo đúng giờ.
% EV: C01; C4A-01; N-21; C10; N-16; N-18; K05; K51.

Trần 120 giây dưới ưu tiên ổn định phục vụ 2.387 lượt, có 2.255 lượt không muộn quá một phút và 2.366 lượt
không muộn quá hai phút trên 3.672 yêu cầu. Đây là một điểm thêm, chưa xác định trần tối ưu. Phép so trực
tiếp luật cộng dồn với luật giới hạn từng lần vẫn chưa có. Bảy lịch sử thỏa điều kiện sàng lọc ở mức 120 giây
trong nhánh ưu tiên chi phí không trần, cũng như một bản ghi chỉ vượt do cộng dồn trong nhánh có trần 30
giây, cho thấy các điều kiện khác nhau của hai luật nhưng chưa cho biết chênh lệch kết quả phục vụ.
% EV: N-06; L1a-02; S15-01 (5d--5e, H7); S15-02; K08.

Ở mật độ cao hơn, trong một bộ mô phỏng đã được kiểm khớp từng khách với RideBound ở cỡ nhỏ và, cho một số luật, ở cỡ Via
trên kịch bản gộp điểm theo khu vực (mục~\ref{muc:c4-cauno}), và trên 60 ngày dữ liệu Via, luật C với ngân sách theo
từng khách phục vụ ngang cách siết đều của Elmi và cộng sự, dạng hàng rào nới, khi chỉ chèn thêm khách (+1,5 khách mỗi
ngày). Khi cho dời khách đang chờ sang xe khác, luật C giảm 16,7 lần đảo hứa mỗi ngày (42\%; khoảng 28--43\% qua các phép
thăm dò khác), đổi lại phục vụ ít hơn 2,8 khách mỗi ngày (0,4\%). So với không gán lại, luật C khi gán lại vẫn có thêm 1,4
lần đảo hứa nhưng phục vụ thêm 4,1 khách mỗi ngày. Ở cùng số khách phục vụ, mỗi khách tránh được đảo hứa đổi lấy khoảng một
khách bị điều phối dời giờ trả muộn quá một phút, nên luật C chỉ đáng dùng nếu một lần bị đảo hứa được coi là khó chịu ít
nhất gần bằng một lần bị dời muộn ($w^* = 0{,}94$, khoảng 0,46--1,61); chưa có số đo nào về cảm nhận của khách. Theo hạn chót
của hàng rào, luật C để nhiều khách đến muộn hơn: 1,78\% so với khoảng 1,1\%. Hiệu đảo hứa giữ dấu khi bộ mô phỏng phá thế
hòa theo cách của RideBound; các hiệu phục vụ nhỏ thì không vững. Mức giảm đảo hứa đến từ chính cơ chế của luật và đảo hứa
được chọn làm chỉ số chính sau thử nghiệm sơ bộ. Đây là một đánh đổi có điều kiện, đo trong một cấu hình thăm dò
(mục~\ref{muc:c4-via-gioihan}), không phải ưu thế chung so với Elmi và cộng sự.
% EV: noi-dung/08c-via-elmi.tex; TABLES-2x3-60days.txt; P5-RESULT.md; breakeven-2026-10-26/BREAKEVEN-RESULT.md; RideBound-Code ADR-085, ADR-086; kiểm toán 23/10 §7-W1, W3, W4; khoảng 28--43\%: báo cáo wf17 §1, §2.3--2.5 (máy công ty).

\section{Vai trò của dự báo và phần còn lại}

Trên các hành trình ưu tiên chi phí có trần đã ghi, tính lại lời hứa có dự báo cho 11 lượt muộn quá một phút
thay vì 72, trong 2.364 lượt phục vụ; lời hứa dài thêm trung bình 7,30 giây. Đối với ưu tiên ổn định có trần,
trên 21 kịch bản ngày thường có dự báo, số lượt muộn giảm từ 68 xuống 2 trong 1.471 lượt phục vụ. Đây là
hai phạm vi khác nhau; dự báo chỉ đổi giờ hứa trong phép tính lại, không đổi việc xe chở ai hay giờ trả thực tế.
% EV: S15-01 (5b); F17; F04; F03; N4B8-01.

Sổ theo từng khách cho biết dự báo giải thích được phần nào của độ lệch. Trong 6.217 giây trễ của 68 khách
ngày thường dưới ưu tiên ổn định, có 82 giây tính cho điều phối, 5.353 giây được hồ sơ lịch sử giải thích và
783 giây còn lại. Chiều tuyết 15/11/2018 cho thấy giới hạn: số khách trễ quá một phút chỉ giảm từ 41 xuống
34 trong 601 khách được phục vụ. Trong 3.445,9 giây trễ của 41 khách đó, 8,6 giây tính cho điều phối, 251,5
giây được hồ sơ lịch sử giải thích và 3.185,9 giây còn lại. Các phần được làm tròn riêng. Phần còn lại không
đồng nghĩa với phần không thể biết trước; sổ là phép phân rã theo mô hình, không xác định lỗi của từng bên.
% EV: F09; N4B8-11; F13; B20; N4B-05; K04; K19; X3-12.

\section{Hạn chế}
\label{muc:c5-hanche}

Các kết quả có tính thăm dò trên 34 kịch bản phụ thuộc nhau thuộc năm ngày đã xem. Mô phỏng có tải thấp,
thời gian dừng đón/trả bằng 0 trong các mẫu đã kiểm và một hệ số tắc đường chung cho toàn thành phố.
Các lượt ưu tiên chi phí dùng bản chạy khác các lượt ưu tiên ổn định; kiểm hành vi trên bốn cặp mẫu chưa
chứng minh tương đương toàn bộ. Chưa có căn cứ suy rộng các chênh lệch thành tác động nhân quả riêng của
bộ chọn, lợi nhuận thực tế, hiệu quả ở tải cao hoặc sự hài lòng của hành khách.

Phần dự báo dùng trung bình lịch sử theo khung 15 phút, lấy điều kiện hiện tại làm mốc; mỗi ngày bật dự báo
chỉ có hai hoặc ba ngày nguồn. Phép kiểm giờ đón vẫn dùng giờ chưa dự báo, trong khi 34 trong 1.471 khách
ngày thường được hiển thị giờ đón đầu vượt hạn đón. Phép thử đưa dự báo trực tiếp vào điều phối có ba lượt
dừng vì kế hoạch không khả thi. Khả năng dự báo và tính nhất quán của phần tích hợp cần được đánh giá thêm.
% EV: I01--I07; B06; B11; B12; S15-01; STAGE1B-PLAN mục 1.1,3; F01--F07; F12.

Khóa luận kế thừa các thành phần đã có: đo và kiểm soát thay đổi lời hứa của Fielbaum và
Alonso-Mora~\cite{fielbaum2020}, mốc tuyến cũ tính lại của Hörl và Zwick~\cite{horl2022}, cùng hạn chót từ lời
hứa đầu của Schulz và Pfeiffer~\cite{schulz2026}. Đối chứng chỉ mượn thành phần hạn chót, chưa tái hiện toàn
bộ các hệ thống đó. Phần thực hiện ở đây là đặc tả, cài đặt, kiểm giới hạn cộng dồn trên phần dời tính cho
điều phối của từng khách, kết hợp lời hứa theo dự báo và sổ theo từng khách, rồi đo các đánh đổi trong phạm vi mô phỏng.

\section{Hướng phát triển}
\label{muc:c5-huongphattrien}

Thứ nhất, cần khảo sát thêm các mức trần dưới cùng một thứ tự chọn phương án, trên dữ liệu phát triển mới.
Trần 30 giây trong khóa luận giới hạn phần dời do điều phối sau khi nhận khách; nó không phải giới hạn toàn bộ
thời gian chờ đón hay đi vòng. Muốn chọn mức triển khai cần biết khách chấp nhận bao nhiêu thay đổi và hãng xe
đánh đổi phục vụ với ổn định thế nào. Không điều chỉnh lại kết quả đã có để chọn một mức có lợi.

Thứ hai, cần mở rộng lịch sử giao thông cùng thành phố từ dữ liệu chuyến taxi theo tháng của
TLC~\cite{tlc}, kiểm chất lượng và độ phủ theo vùng, ngày, giờ trước khi dùng. Phải tách thời gian học, chọn
cấu hình và kiểm tra trên ngày chưa xem; mô hình phức tạp hơn chỉ được giữ nếu cải thiện ngoài tập học.
Ngày có sự kiện cần được đánh giá riêng, chỉ dùng tín hiệu thực sự có ở thời điểm hứa. Dữ liệu mô phỏng
bổ sung hữu ích cho thử tình huống khó, nhưng không thay bằng chứng dự báo trên giao thông thực.

Thứ ba, cần sửa phép kiểm giờ đón để phù hợp với lời hứa theo dự báo, đồng thời kiểm dự báo sai, dữ liệu
đến chậm và trường hợp thiếu lịch sử. Việc sửa cần có phép thử mô phỏng đầy đủ ở cấu hình mới; kết quả
tính lại lời hứa trên hành trình cũ không xác nhận thay đổi trong quyết định nhận và ghép khách.

Thứ tư, phép so luật cộng dồn với luật giới hạn từng lần vẫn là hướng kiểm tra riêng, cần đủ tình huống
phân biệt hai luật và kế hoạch phù hợp. Cũng cần thử hạn chót nới với nhiều mức trần hơn, mở rộng tải,
tắc đường riêng trên từng đoạn và mô hình khách hủy chuyến. Các hướng này giữ nguyên kết quả và giới hạn
của những phép thử đã trình bày, thay vì coi các phần chưa đo là ưu thế đã có.
